% =====================================================
% 题型：锐角三角形面积与余弦定理填空题 (q_tjw_20260606_22_acute_area_side.tex)
% =====================================================
% =====================================================
% 难度：★★ (中档)
% =====================================================
\newcommand{\qTriangleObtuseAreaSideStem}{%
  若锐角 \(\triangle ABC\) 的面积为 \(10\sqrt{3}\)，且 \(AB = 5\)，\(AC = 8\)，
  则 \(BC = \underline{\hspace{2.5cm}}\)。%
}

\newcommand{\qTriangleObtuseAreaSideAnswer}{%
  \(7\)%
}

\newcommand{\qTriangleObtuseAreaSideSolution}{%
  由面积公式得
  \[
    S_{\triangle ABC}
    = \frac{1}{2} \cdot AB \cdot AC \sin A
    = \frac{1}{2} \cdot 5 \cdot 8 \sin A
    = 20\sin A.
  \]
  因为面积为 \(10\sqrt{3}\)，所以
  \[
    20\sin A = 10\sqrt{3},
    \qquad
    \sin A = \frac{\sqrt{3}}{2}.
  \]

  又 \(\triangle ABC\) 为锐角三角形，且 \(AB=5, AC=8\) 夹角为 \(A\)。
  若 \(A=60^\circ\)，则
  \[
    BC^2 = 5^2+8^2-2\cdot 5\cdot 8\cos60^\circ = 49,
  \]
  此时三边为 \(5,7,8\)，且 \(8^2<5^2+7^2\)，符合锐角三角形条件。

  若 \(A=120^\circ\)，则 \(BC^2=129\)，此时最大边为 \(BC\)，但
  \[
    129 > 5^2+8^2=89,
  \]
  会使角 \(A\) 为钝角，不符合题意。

  因此
  \[
    BC = 7.
  \]
}

\newcommand{\qTriangleObtuseAreaSideDiagnosis}{%
  （留白待收集学生作答后补充）%
}

\newcommand{\qTriangleObtuseAreaSideTeachingNotes}{%
  精讲候选。表面是面积公式加余弦定理，实质要讨论 \(\sin A=\frac{\sqrt3}{2}\) 对应 \(A=60^\circ\) 或 \(120^\circ\)。本题用“锐角三角形”排除 \(A=120^\circ\) 的分支。%
}


% =====================================================
% 别名兼容：旧课次宏定义
% =====================================================
\newcommand{\qTJWZeroSixTwentyTwoStem}{\qTriangleObtuseAreaSideStem}
\newcommand{\qTJWZeroSixTwentyTwoAnswer}{\qTriangleObtuseAreaSideAnswer}
\newcommand{\qTJWZeroSixTwentyTwoSolution}{\qTriangleObtuseAreaSideSolution}
\newcommand{\qTJWZeroSixTwentyTwoDiagnosis}{\qTriangleObtuseAreaSideDiagnosis}
\newcommand{\qTJWZeroSixTwentyTwoTeachingNotes}{\qTriangleObtuseAreaSideTeachingNotes}
